\ifdefined\standalone
\documentclass[12pt]{article}
\usepackage{graphicx}
\usepackage{float}
\usepackage[a4paper, margin=1in]{geometry}
\usepackage[utf8]{inputenc}
\usepackage{textcomp}
\usepackage{amsmath}
\usepackage{booktabs}
\usepackage{tikz}
\usepackage{enumitem}
\usetikzlibrary{decorations.pathmorphing,patterns}
\usepackage[hidelinks]{hyperref}

\begin{document}
\fi


\section{Convergence Of The Thermal Solver}

This verification case evaluates the temporal and spatial convergence of the one-dimensional thermal solver.  
The configuration corresponds to the case presented in Section~\ref{case:rad_1d}.  
A slab of thickness \(4~\mathrm{cm}\) is considered, with a prescribed heat flux applied at the upper boundary.  


\subsection*{Temporal convergence}

The temporal convergence is examined by varying the time step from \(0.01~\mathrm{s}\) to \(100~\mathrm{s}\), while maintaining a uniform spatial discretization of \(0.1~\mathrm{mm}\).  
Although the implicit formulation ensures stability for all tested time steps, significant deviations from the analytical solution are expected when coarse temporal resolutions are used.

The temperature at a depth of \(8~\mathrm{mm}\) is monitored over time.  
Figure~\ref{fig:Temperature_Temporal_convergence} presents the computed temperature histories, and Figure~\ref{fig:L2_error_Temporal_thermal_solver} shows the associated \(L^2\) error (Eq. \ref{eq:L2norm}) relative to the analytical solution. A linear slope equal to 1 is obtained in the log–log error plot, which is expected since the thermal solver is first order accurate in time.

\begin{figure}[H]
\centering
\begin{minipage}{0.48\textwidth}
    \centering
    \includegraphics[width=\textwidth]{Temporal/Temperature_Temporal_convergence.png}
    \caption{Temperature at 8 mm depth for various time steps.}
    \label{fig:Temperature_Temporal_convergence}
\end{minipage}
\hfill
\begin{minipage}{0.48\textwidth}
    \centering
    \includegraphics[width=\textwidth]{Temporal/L2_error_Temporal_thermal_solver.png}
    \caption{\(L^2\) temporal error as a function of time step.}
    \label{fig:L2_error_Temporal_thermal_solver}
\end{minipage}
\end{figure}


\subsection*{Spatial convergence}

The spatial convergence is evaluated by varying the grid spacing from \(10~\mathrm{mm}\) (five cells through the slab) down to \(0.001~\mathrm{mm}\).  
All simulations are run using a fixed time step of \(0.01~\mathrm{s}\).  
Unconditional stability is preserved regardless of the grid refinement, although coarse meshes are expected to deviate from the analytical solution.

The temperature at \(8~\mathrm{mm}\) depth is again used to compute the \(L^2\) error. Figure~\ref{fig:Temperature_Spatial_convergence} presents the computed temperature histories. Figure~\ref{fig:L2_error_Spatial_thermal_solver} presents the spatial error as a function of grid spacing. A linear slope equal to 2 is obtained in the log–log error plot, which is expected since the thermal solver is second order accurate in space.

\begin{figure}[H]
\centering
\begin{minipage}{0.47\textwidth}
    \centering
    \includegraphics[width=\textwidth]{Spatial/Temperature_Spatial_convergence.png}
    \caption{Temperature at 8 mm depth for various time steps.}
    \label{fig:Temperature_Spatial_convergence}
\end{minipage}
\hfill
\begin{minipage}{0.47\textwidth}
    \centering
    \includegraphics[width=\textwidth]{Spatial/L2_error_Spatial_thermal_solver.png}
    \caption{\(L^2\) spatial error as a function of grid spacing.}
    \label{fig:L2_error_Spatial_thermal_solver}
\end{minipage}
\end{figure}


\subsection*{Computational cost}

The computational cost of the simulations is reported in Figures~\ref{fig:CPU_Time_Temporal_thermal_solver} and \ref{fig:CPU_Time_Spatial_thermal_solver}.  
A proportional increase in CPU time is expected as the spatial or temporal resolution is refined.  
For coarse discretizations, however, initialization and overhead operations may dominate the total runtime, and meaningful scaling behaviour is observed only once the discretization becomes sufficiently fine.

\begin{figure}[H]
\centering
\begin{minipage}{0.48\textwidth}
    \centering
    \includegraphics[width=\textwidth]{Temporal/CPU_Time_Temporal_thermal_solver.png}
    \caption{CPU time for temporal convergence simulations.}
    \label{fig:CPU_Time_Temporal_thermal_solver}
\end{minipage}
\hfill
\begin{minipage}{0.48\textwidth}
    \centering
    \includegraphics[width=\textwidth]{Spatial/CPU_Time_Spatial_thermal_solver.png}
    \caption{CPU time for spatial convergence simulations.}
    \label{fig:CPU_Time_Spatial_thermal_solver}
\end{minipage}
\end{figure}




\ifdefined\standalone
\end{document}
\fi


